% SPDX-License-Identifier: LPPL-1.3c
% Copyright (C) 2026 Christophe Jorssen
% Author and Current Maintainer: Christophe Jorssen <christophe.jorssen@gmail.com>
% LPPL maintenance status: maintained
% This file is part of luafemm. See LICENSE and LICENSES.md for its terms.
\documentclass[a4paper,11pt]{article}
\usepackage[margin=19mm]{geometry}
\usepackage{fontspec}
\usepackage{amsmath}
\usepackage{luafemm}
\pagestyle{empty}
\definecolor{alnico}{HTML}{B75046}
\definecolor{flux}{HTML}{147B83}
\begin{document}
\noindent{\LARGE\bfseries A U-shaped alnico and soft-iron magnet}\par
\smallskip
\noindent\textsf{luafemm \luafemmversion{} · nonlinear permanent magnet · zero imposed current}

\medskip
\noindent An alnico crosspiece magnetized to the left drives two pole pieces
and a soft-iron armature. The flux crosses two $1\,\mathrm{mm}$ gaps,
with local mesh refinement.

\tikzset{
  femm/materials/iron/.style={bh={0/0,.5/80,1/180,1.3/350,
    1.5/900,1.65/2500,1.8/7000,2/20000}},
  % FEMM matlib.dat: Cast Alnico 5 (LNG37).
  % SHIFTED H(B) table: the solver subtracts Hc along the magnetization.
  femm/materials/alnico/.style={coercivity=48600,
    bh={0/0,.2781/1600,.6383/4100,.8234/6300,.9214/8600,
      1.013/13400,1.064/18200,1.106/24300,1.149/33000,
      1.184/42530,1.204/48600}},
  iron/.style={femm/region={material=iron},draw=black!65,fill=black!9},
  magnet/.style={femm/region={material=alnico,magnetization angle=180},
    draw=alnico!80!black,fill=alnico!25}
}
\begin{center}
\begin{tikzpicture}[femm/problem={xmin=-100,xmax=100,ymin=-85,ymax=100,
  mesh size=2,mesher=delaunay}]
  \filldraw[iron] (-40,-30)--(40,-30)--(40,30)--(25,30)
    --(25,-15)--(-25,-15)--(-25,30)--(-40,30)--cycle;
  \filldraw[magnet] (-25,-30) rectangle (25,-15);
  \filldraw[iron] (-40,31) rectangle (40,46);
  % Two air regions also prescribe local mesh refinement.
  \path[femm/region={mesh size=.5}] (-40,30) rectangle (-25,31);
  \path[femm/region={mesh size=.5}] (25,30) rectangle (40,31);
  \begin{scope}
    \clip (-61,-45) rectangle (62,64);
    \pic[femm/lines=29,femm/arrow spacing=26mm,
      femm/field style={flux}] {femm field};
  \end{scope}
  \node[fill=white,inner sep=2pt,font=\small] at (0,39) {Soft-iron armature};
  \node[fill=white,inner sep=2pt,font=\small] at (0,4) {Air: $J_z=0$ everywhere};
  \node[fill=white,inner sep=2pt,text=flux] at (-32.5,27) {N};
  \node[fill=white,inner sep=2pt,text=flux] at (32.5,27) {S};
  \node[fill=alnico!25,inner sep=2pt,font=\small] at (0,-20) {Alnico 5 (LNG37)};
  \draw[white,line width=2.5pt] (15,-25)--(-15,-25);
  \draw[-Stealth,line width=1pt,alnico!80!black] (15,-25)--(-15,-25)
    node[midway,below,fill=alnico!25,inner sep=1pt,font=\scriptsize] {magnetization direction};
  \draw[thin] (40,30)--(50,30) (40,31)--(50,31);
  \draw[thin,-Stealth] (48,25)--(48,30);
  \draw[thin,-Stealth] (48,36)--(48,31);
  \node[anchor=west,fill=white,font=\small,inner sep=1pt] at (50,30.5) {$e=1\,\mathrm{mm}$};
  \draw[thin] (-25,-30)--(-25,-38) (25,-30)--(25,-38);
  \draw[Stealth-Stealth,thin] (-25,-36)--(25,-36)
    node[midway,fill=white,font=\small,inner sep=2pt] {$50\,\mathrm{mm}$};
  \node[fill=white,font=\small,text=flux,inner sep=2pt] at (0,58) {$\mathbf B$ contours, oriented from N to S in air};
\end{tikzpicture}
\end{center}

\noindent\textbf{Computed result.}
At the gap centres, $|\mathbf B|=\femmB{-32.5}{30.5}\,\mathrm T$
on the left and $\femmB{32.5}{30.5}\,\mathrm T$ on the right.
The mesh contains \femmstat{nodes} nodes and \femmstat{triangles} triangles;
Newton converges in \femmstat{newton} iterations.

\medskip
\noindent\textbf{Magnet law.}
The model uses $\mathbf H=\nu(|\mathbf B|)\mathbf B-H_c\mathbf m$,
where $\mathbf m$ is the prescribed direction and $H_c=48.6\,\mathrm{kA\,m^{-1}}$.
The shifted FEMM material table gives a remanence of
$B_r=1.204\,\mathrm T$. It describes a demagnetization branch,
without hysteresis history or irreversible demagnetization.

\medskip
\noindent\textbf{TikZ declaration.}
\begin{verbatim}
femm/materials/alnico/.style={coercivity=48600,bh={...}}
femm/region={material=alnico,magnetization angle=180}
femm/region={material=air,mesh size=.5}
\end{verbatim}
\noindent\small Material data and model limits: \texttt{docs/permanent-magnets.md}.
The mesher uses filtered geometric predicates with exact Lua fallbacks.
\end{document}
